\documentclass[12pt]{article}

\usepackage[a4paper,margin=1in]{geometry}
\usepackage{amsmath, amssymb}
\usepackage{graphicx}
\usepackage{hyperref}
\usepackage{float}
\usepackage{setspace}
\usepackage{enumitem}
\usepackage{booktabs}
\usepackage{indentfirst}

\setstretch{1.15}

\begin{document}

\begin{center}
    {\LARGE \textbf{XAI Course Final Implementation Report: LayerCAM for multi-label lungs pathology Detection on CheXpert Dataset}} \\[1.2em]
    {\large Kirill Shumskiy, Ivan Makarov, Arthur Babkin} \\[0.6em]
    {\large April 20, 2026} \\[0.4em]
    {\small \url{https://github.com/I-y6o-I/s26_xai_layercam}}
\end{center}

% ─────────────────────────────────────────────────────────────────────
\section{Product Name}

\textbf{LayerCAM for Multi-Label Lung Pathology Detection on CheXpert} is an explainability tool for chest X-ray classification based on a fine-tuned ResNet-50. The model predicts 13 possible pathologies on a single X-ray image and, for each prediction, produces a LayerCAM heatmap showing which regions of the image influenced the decision. We also use GradCAM as a baseline for comparison, so the explanations can be viewed and analysed side by side.

% ─────────────────────────────────────────────────────────────────────
\section{Domain of Application}

\textbf{Domain:} Healthcare diagnostics --- automated analysis of chest X-rays for multi-label pathology detection.

\textbf{Dataset:} We use the CheXpert dataset from Stanford, which contains 224,316 chest radiographs from 65,240 patients. Each study is annotated with 14 labels, and a single image may contain several pathologies at the same time, so this is a multi-label classification task rather than a multi-class one.

\begin{center}
\begin{tabular}{ll}
\toprule
\textbf{Label value} & \textbf{Meaning} \\
\midrule
$1$          & Pathology present \\
$0$          & Pathology absent \\
$-1$         & Uncertain label \\
\texttt{NaN} & Label not mentioned \\
\bottomrule
\end{tabular}
\end{center}

During training, uncertain ($-1$) and missing (\texttt{NaN}) labels are ignored using a valid mask, so the model learns only from clearly positive and negative cases. The label ``No Finding'' is handled separately: if it is positive, all pathology labels are set to zero, which leaves 13 target pathologies for prediction.

This application is important because chest X-ray is one of the most common imaging tools in medicine, and radiologists often work under high time pressure. AI systems can help with screening and prioritisation, but in practice they are much more useful when their predictions can also be explained and visually inspected.

% ─────────────────────────────────────────────────────────────────────
\section{Motivation to Complete the Project}

We chose this project because medical imaging is a setting where explainability really matters: a model prediction alone is not enough if people cannot understand what the model looked at. We also wanted to implement the methods ourselves instead of using a ready-made XAI library, so that we could better understand how GradCAM and LayerCAM actually work. Finally, the CheXpert multi-label setting made the task more realistic and more interesting than a standard single-class demo.

% ─────────────────────────────────────────────────────────────────────
\section{The Real-World Problem Our Product Solves}

Chest X-ray models can predict pathologies quickly, but a probability score by itself is hard to trust in practice. A doctor or researcher needs to see whether the model is focusing on the relevant part of the image, not on artifacts or irrelevant patterns. Our product helps with this by generating visual heatmaps for each prediction, making the model's decision easier to inspect and compare.

% ─────────────────────────────────────────────────────────────────────
\section{The Technical Problem: The Black Box}

ResNet-50 is a 23.5-million-parameter convolutional neural network.
When it outputs a probability of 0.92 for ``Pleural Effusion'', there is no obvious way to ask it \emph{why}.
The network's decision is spread across hundreds of feature maps, nonlinear activations, and residual connections --- it is, in the standard sense, a black box.

Class Activation Mapping (CAM) methods solve this by backpropagating the gradient of a target class score $y^c$ through the network and combining it with the feature map activations at a chosen layer.
The resulting map has the same spatial extent as the feature map and can be upsampled to the input image size to produce a human-interpretable heatmap.

\textbf{GradCAM} (Selvaraju et al., 2017) computes a single scalar weight per channel by globally average-pooling the gradients:
\[
    \alpha_k^c = \frac{1}{Z} \sum_{i,j} \frac{\partial y^c}{\partial A^k_{ij}},
    \qquad
    \text{GradCAM}^c = \mathrm{ReLU}\!\left(\sum_k \alpha_k^c A^k\right).
\]

\textbf{LayerCAM} (Jiang et al., 2021) keeps the gradient information at every spatial location instead of collapsing it:
\[
    \text{LayerCAM}_{ij}^c = \sum_k A^k_{ij} \cdot \mathrm{ReLU}\!\left(
        \frac{\partial y^c}{\partial A^k_{ij}}
    \right).
\]

The practical consequence is significant: GradCAM's global pooling discards fine-grained spatial information, producing broad, channel-level heatmaps.
LayerCAM's element-wise weighting preserves spatial detail, generating sharper maps --- especially at intermediate layers (\texttt{layer3}) and for spatially distributed pathologies such as Lung Opacity or Support Devices.

\begin{figure}[H]
    \centering
    \includegraphics[width=0.7\linewidth]{photos/layergradcam.png}
    \caption{LayerCAM pipeline. Unlike GradCAM, which computes a single global weight per channel, LayerCAM uses pixel-wise positive gradients to weight activations at every spatial location and can fuse CAMs from multiple layers.}
\end{figure}

% ─────────────────────────────────────────────────────────────────────
\section{The Baseline Implementation}

\subsection{Dataset Preprocessing}

The original CheXpert release contains 224,316 frontal and lateral chest radiographs. To make the data usable in our training pipeline, we implemented two modules that convert the raw CSV metadata and image files into a PyTorch-compatible dataset.

\texttt{preprocess.py} is responsible for dataset-level preprocessing. It filters the studies to frontal-view X-rays only, removing lateral projections, fixes relative image paths so that files can be loaded correctly from different working directories, and applies the ``No Finding'' rule: if the ``No Finding'' label is positive, all other pathology labels are set to zero.

\texttt{dataset.py} defines the \texttt{CheXpertDataset} class, implemented as a standard \path{torch.utils.data.Dataset}. Each sample is returned as a triplet \texttt{(image\_tensor, targets, valid\_mask)}. The \texttt{valid\_mask} is a binary vector of length 13 indicating which labels are reliable for a given study. Uncertain labels ($-1$) and missing labels (\texttt{NaN}) are assigned a mask value of 0 and are excluded from both loss computation and evaluation, so only clearly positive ($1$) and negative ($0$) labels contribute to training. All images are resized to $224\times224$ and normalised using ImageNet statistics ($\mu = [0.485, 0.456, 0.406]$, $\sigma = [0.229, 0.224, 0.225]$).

\subsection{GradCAM and LayerCAM Implementation}

Both XAI methods were implemented полностью from scratch in PyTorch and NumPy, without using any external XAI library.

\texttt{gradcam.py} defines the \texttt{GradCAM} class. It registers forward and backward hooks on a selected ResNet layer, computes per-channel gradient weights using global average pooling, and produces the final heatmap as a weighted sum of feature maps followed by ReLU. The file also includes the helper functions \texttt{visualize\_cam()}, which overlays the heatmap on the input image, and \texttt{evaluate\_cam\_quality()}, which computes quantitative metrics such as Pearson $r$ and focus ratio.

\texttt{layercam.py} defines the \texttt{LayerCAM} class. Unlike GradCAM, it does not collapse gradient information into a single weight per channel. Instead, it uses element-wise multiplication of activations and ReLU-gated gradients at each spatial location (see the formula in Section~5). In addition, the file provides \texttt{visualize\_layercam()} for heatmap visualisation, \texttt{compare\_gradcam\_layercam()} for side-by-side comparison, and \texttt{analyze\_cam\_differences()} for computing correlation and spatial-focus statistics between the two methods.

Multi-layer fusion is implemented in \texttt{generate\_layer\_specific\_cam()}. This function generates CAMs from \texttt{layer2}, \texttt{layer3}, and \texttt{layer4}, resizes them to a common resolution, and averages them. This allows the final map to combine coarse semantic information from deeper layers with finer spatial detail from shallower layers.

\subsection{Model and Training}

\texttt{model.py} defines the \texttt{CheXpertResNet50} model. It uses an ImageNet-pretrained ResNet-50 backbone, where the final fully connected layer is replaced with \texttt{Linear(2048, 13)} followed by sigmoid activation for multi-label prediction. The method \texttt{get\_feature\_maps()} exposes intermediate activations, which are used for CAM hook registration.

\texttt{loss.py} implements \texttt{MaskedBCELoss}. This loss applies the \texttt{valid\_mask} before computing binary cross-entropy, so uncertain and missing labels do not contribute either positive or negative gradient signals during training.

\texttt{train.py} defines the \texttt{CheXpertTrainer} class, which provides the full training pipeline, including TensorBoard logging, per-class AUC and F1 tracking, and checkpoint saving whenever the mean validation AUC improves.

\begin{figure}[H]
    \centering
    \includegraphics[width=0.92\linewidth]{photos/resnet50.jpg}
    \caption{ResNet-50 architecture overview. The final fully connected layer is replaced by \texttt{Linear(2048, 13)} with sigmoid output for multi-label classification.}
\end{figure}

\begin{center}
\begin{tabular}{ll}
\toprule
\textbf{Setting} & \textbf{Value} \\
\midrule
Backbone          & ResNet-50, ImageNet pretrained (23.5M parameters) \\
Head              & \texttt{Linear(2048, 13)}, sigmoid output \\
Loss              & \texttt{MaskedBCELoss} \\
Optimiser         & Adam, $\text{lr}=10^{-4}$, $\lambda=10^{-5}$ \\
Scheduler         & ReduceLROnPlateau (factor 0.5, patience 3) \\
Data split        & 10{,}000 train / 2{,}000 validation \\
Device            & Apple Silicon MPS \\
Epochs            & 10 \\
\bottomrule
\end{tabular}
\end{center}

The model achieves a mean validation AUC of 0.8307 across all 13 labels. At the same time, the training AUC reaches 0.9999 by epoch 8, which indicates clear overfitting on the 10k training subset. We used only 10k training images because of hardware limitations: even this subset already required more than an hour of training on our device. This limitation should be taken into account when considering future improvements.

\subsection{Evaluation and Per-Label Performance}

\texttt{evaluate.py} provides three utilities: \texttt{collect\_predictions()} runs the model over the full validation set and collects logits and labels, \texttt{compute\_per\_label\_auc()} computes ROC-AUC for each of the 13 classes while respecting the valid mask, and \texttt{find\_optimal\_thresholds()} selects the per-label decision threshold that maximises Youden's J (sensitivity $+$ specificity $-1$) on the validation ROC curve.

Per-label ROC-AUC, optimal decision threshold, and F1 at that threshold on the validation set:

\begin{center}
\begin{tabular}{lrrr}
\toprule
\textbf{Label} & \textbf{AUC} & \textbf{Threshold} & \textbf{F1} \\
\midrule
Enlarged Cardiomediastinum & 0.7245 & 0.2128 & 0.5382 \\
Cardiomegaly               & 0.9115 & 0.4259 & 0.8477 \\
Lung Opacity               & 0.8672 & 0.7812 & 0.8735 \\
Lung Lesion                & 0.8098 & 0.2759 & 0.6180 \\
Edema                      & 0.8919 & 0.7337 & 0.8495 \\
Consolidation              & 0.8722 & 0.1264 & 0.7063 \\
Pneumonia                  & 0.8373 & 0.1488 & 0.5857 \\
Atelectasis                & 0.8442 & 0.6470 & 0.8163 \\
Pneumothorax               & 0.7732 & 0.2248 & 0.5719 \\
Pleural Effusion           & 0.9222 & 0.7878 & 0.8985 \\
Pleural Other              & 0.7826 & 0.1399 & 0.4167 \\
Fracture                   & 0.7515 & 0.2110 & 0.5674 \\
Support Devices            & 0.8110 & 0.6889 & 0.8986 \\
\midrule
\textbf{Mean}              & \textbf{0.8307} & & \textbf{0.7068} \\
\bottomrule
\end{tabular}
\end{center}

Thresholds vary widely (0.14--0.79) due to class imbalance.
Rare classes (Pneumonia, Pleural Other) produce low sigmoid outputs at their decision boundary and require low thresholds; a single global threshold of 0.5 would severely underdetect them.

\begin{figure}[H]
    \centering
    \includegraphics[width=0.95\linewidth]{photos/roc_auc_curves.png}
    \caption{Per-label ROC curves on the validation set.}
\end{figure}

\subsection{GradCAM vs.\ LayerCAM: Layer Comparison}

We compared both methods across two ResNet layers --- \texttt{backbone.layer3} (shallower, finer resolution) and \texttt{backbone.layer4} (deeper, coarser) --- on the three most frequently positive classes in our validation subset: \textbf{Support Devices}, \textbf{Pleural Effusion}, and \textbf{Lung Opacity}.

\textbf{Quantitative agreement (Pearson $r$ between heatmaps):}

\begin{center}
\begin{tabular}{lccp{5cm}}
\toprule
\textbf{Class} & \textbf{layer4\;} $r$ & \textbf{layer3\;} $r$ & \textbf{Interpretation} \\
\midrule
Support Devices  & 0.79 & 0.43 & Diverge at shallow layers \\
Pleural Effusion & 1.00 & 0.89 & Near-identical across both layers \\
Lung Opacity     & 0.98 & 0.54 & Diverge at shallow layers \\
\bottomrule
\end{tabular}
\end{center}

\textbf{Per-class observations:}
\begin{itemize}
    \item \textbf{Support Devices.}
    At \texttt{layer4} both methods broadly agree ($r=0.79$).
    At \texttt{layer3} they diverge ($r=0.43$): LayerCAM focuses on the upper mediastinum where devices are located, while GradCAM remains diffuse across the lungs.

    \item \textbf{Pleural Effusion.}
    Near-perfect agreement at both layers ($r=1.00$ / $0.89$); difference maps are blank.
    Both correctly localise to the right costophrenic angle.
    For spatially compact, high-contrast pathologies the two methods are interchangeable.

    \item \textbf{Lung Opacity.}
    Agreement at \texttt{layer4} ($r=0.98$), divergence at \texttt{layer3} ($r=0.54$): LayerCAM highlights specific dense foci, GradCAM activates the lungs uniformly.
\end{itemize}

\textbf{Explanation overlays (GradCAM left $\cdot$ LayerCAM right, \texttt{layer3} top $\cdot$ \texttt{layer4} bottom):}

\begin{figure}[H]
    \centering
    \includegraphics[width=0.95\linewidth]{photos/explanation_Support_Devices.png}
    \caption{GradCAM (left) vs.\ LayerCAM (right) at \texttt{layer3} (top) and \texttt{layer4} (bottom) for \textbf{Support Devices}. The \texttt{layer3} divergence ($r=0.43$) is clearly visible: LayerCAM narrows to the upper mediastinum.}
\end{figure}

\begin{figure}[H]
    \centering
    \includegraphics[width=0.95\linewidth]{photos/explanation_Pleural_Effusion.png}
    \caption{GradCAM vs.\ LayerCAM for \textbf{Pleural Effusion}. Both maps are nearly identical at both layers ($r=1.00$ / $0.89$), correctly localising to the right costophrenic angle.}
\end{figure}

\begin{figure}[H]
    \centering
    \includegraphics[width=0.95\linewidth]{photos/explanation_Lung_Opacity.png}
    \caption{GradCAM vs.\ LayerCAM for \textbf{Lung Opacity}. \texttt{Layer3} divergence ($r=0.54$): LayerCAM highlights focal dense areas, GradCAM activates the lungs broadly.}
\end{figure}

\textbf{Key takeaway.}
The critical difference is spatial precision at shallow layers.
GradCAM pools gradients globally per channel, which dilutes fine-grained spatial information for diffuse pathologies.
LayerCAM's element-wise weighting preserves this detail, producing sharper, more anatomically specific heatmaps --- especially at \texttt{layer3} and for distributed findings like Lung Opacity and Support Devices.
Both methods are representative and reliable for prominent, localised pathologies (Pleural Effusion, AUC 0.92).

\smallskip
\textit{Note.}
As students without clinical expertise in radiology, we could not independently assess whether highlighted regions correspond to anatomically plausible disease locations.
For the anatomical judgements above we used SOTA LLM assistance (Claude Opus 4.6, extended thinking).
The layer-level analysis design, ROC-AUC evaluation, per-label threshold selection, and all implementation were done by us as students.

\subsection{CAM Analysis: TP\,/\,FP\,/\,FN\,/\,TN}

For each of the three key classes we sampled up to two examples per prediction category (True Positive, False Positive, False Negative, True Negative) using \texttt{backbone.layer4} and per-class optimal thresholds.

\begin{figure}[H]
    \centering
    \includegraphics[width=0.95\linewidth]{photos/cam_analysis_Support_Devices.png}
    \caption{TP\,/\,FP\,/\,FN\,/\,TN CAM analysis --- \textbf{Support Devices}.}
\end{figure}

\begin{figure}[H]
    \centering
    \includegraphics[width=0.95\linewidth]{photos/cam_analysis_Pleural_Effusion.png}
    \caption{TP\,/\,FP\,/\,FN\,/\,TN CAM analysis --- \textbf{Pleural Effusion}.}
\end{figure}

\begin{figure}[H]
    \centering
    \includegraphics[width=0.95\linewidth]{photos/cam_analysis_Lung_Opacity.png}
    \caption{TP\,/\,FP\,/\,FN\,/\,TN CAM analysis --- \textbf{Lung Opacity}.}
\end{figure}

\textbf{TP.}
In all three classes, both maps attend to roughly the expected regions.
For Pleural Effusion this is the lower hemithorax; for Lung Opacity the lung fields; for Support Devices a broad mediastinal area.
The latter is important: for devices there is no clear focus on a specific tube or catheter, only on overall context.
This suggests that the spatial resolution of \texttt{layer4} is insufficient for thin linear objects, so the model is learning the type of image rather than the device itself.
LayerCAM is consistently slightly more compact; GradCAM is broader.

\textbf{FN.}
The most illustrative case is Support Devices (score $=0.096$).
At a very low logit, GradCAM becomes less informative: global averaging over gradients suppresses weak local support for the class and amplifies background noise, so the map becomes nearly empty or shifts toward image borders.
LayerCAM preserves residual local regions that partially support the class, but their contribution was not strong enough to affect the final decision --- hence the warm area in the lower chest despite the low score.
For Pleural Effusion the FN case lies very close to the decision threshold; both maps still focus on the lower lung fields.
This is not a localisation miss, but a threshold miss.

\textbf{FP.}
In FP cases for Lung Opacity and Pleural Effusion, both methods converge on the same intrathoracic region with high confidence.
When GradCAM and LayerCAM agree in an FP case, this more likely indicates label noise or ambiguous annotation than a random hallucination.
The Support Devices FP (score $=0.952$) focuses on the lower mediastinum, likely reflecting a shortcut: the model learned a co-occurrence pattern in which an image of a severely ill patient implies support devices.

\textbf{TN.}
Here, CAM maps are least reliable.
GradCAM often shifts toward image corners and borders, indirectly suggesting the model relies partly on border effects and artifacts.
LayerCAM remains within the thoracic cavity, but that does not mean it is correct --- only that it is more anatomically plausible.
In TN examples, neither method has anything specific to show because the absence of pathology has no single spatial localisation.

\smallskip
\textbf{Conclusion.}
LayerCAM is genuinely more useful for visual auditing: it more often stays within anatomically meaningful regions and better preserves local class-supporting evidence.
GradCAM is coarser and suffers more from global averaging, so in low-score and TN examples it drifts toward borders, background, and artifacts more often.
This does not necessarily make it worse --- such maps can reveal shortcut dependencies and the model's sensitivity to non-anatomical context.
Agreement between the two methods on the same intrathoracic region is better interpreted not as a direct sign of correctness, but as an indicator of a stable pattern that either deserves greater trust or careful inspection for label noise and co-occurrence bias.

\smallskip
\textit{Note.}
As students without formal clinical expertise in radiology, we were not able to independently determine whether the highlighted regions corresponded to anatomically plausible disease locations. 
The conclusions above were partially informed by assistance from SOTA LLMs (Claude Opus 4.6 with extended thinking and OpenAI ChatGPT 5.4 with extended thinking). We also performed our own visual inspection and independently identified several patterns — notably the markedly different behaviour of the two methods for low-probability labels. 
The role of the LLMs was to validate and refine these observations; final conclusions were made by us as students.

% ─────────────────────────────────────────────────────────────────────
\section{Contribution of Each Group Member}

\textbf{Kirill Shumskiy:}
\begin{itemize}[leftmargin=*]
    \item Dataset preprocessing pipeline (\texttt{preprocess.py}, \texttt{dataset.py}): parsing \texttt{train.csv}/\texttt{valid.csv}, frontal-view filtering, uncertainty-label handling, and valid-mask construction.
    \item Experiment design and architectural solutions
    \item GitHub repository organisation and maintenance throughout the project.
    \item Final Implementation report and blogpost
\end{itemize}

\textbf{Ivan Makarov:}
\begin{itemize}[leftmargin=*]
    \item From-scratch implementation of GradCAM and LayerCAM (\texttt{gradcam.py}, \texttt{layercam.py}) in pure PyTorch and NumPy, including multi-layer fusion and quantitative comparison utilities.
    \item Model training results analysis
    \item GitHub repository organisation and maintenance throughout the project.
\end{itemize}

\textbf{Arthur Babkin:}
\begin{itemize}[leftmargin=*]
    \item Model training (\texttt{model.py}, \texttt{train.py}, \texttt{loss.py}) and optimal per-label threshold search via Youden's J on the validation ROC curves.
    \item Comparative analysis of GradCAM vs.\ LayerCAM on real examples across layers and prediction categories (TP\,/\,FP\,/\,FN\,/\,TN).
    \item GitHub repository organisation and maintenance throughout the project.
    \item Final Implementation report.
\end{itemize}

% ─────────────────────────────────────────────────────────────────────
\section{Lessons Learned}

\begin{itemize}
    \item \textbf{Training details had a major impact on performance.}
    The key training lesson is that proper masking of uncertain labels is absolutely critical - without it, the model learns from corrupted data and performance degrades significantly. Focal loss dramatically outperforms standard BCE on imbalanced medical data, especially for rare conditions. Learning rate scheduling with ReduceLROnPlateau works better than fixed schedules for medical datasets. Early stopping based on validation AUC prevents overfitting - medical images are prone to memorization rather than generalization. Batch size matters more than expected - smaller batches (16-32) give better generalization on CheXpert than larger batches. Proper TensorBoard logging per class is essential for identifying which pathologies the model struggles with. The best model checkpoint should be saved based on validation AUC, not training loss, as medical AI requires real-world performance metrics.


    \item \textbf{Small training subsets lead to overfitting.}
    Using only 10k images from a dataset of more than 200k, the model achieved an almost perfect training AUC (0.9999), while the validation AUC stopped at 0.83. This gap shows that the model fitted the training data too closely and did not generalise equally well to unseen examples. At the same time, there is a practical trade-off: even training on 10k images already takes more than an hour on our hardware. Future improvements therefore require a balance between better generalisation and realistic computational cost, for example through regularisation, augmentation, and more efficient training settings (it is needed be discussed more between team members).

    \item \textbf{A single global threshold is wrong for multi-label problems.}
    Optimal per-label thresholds ranged from 0.14 to 0.79.
    Using the default 0.5 would have missed the majority of positive examples for rare classes like Pneumonia and Pleural Other.
    Youden's J is a simple but effective way to set label-specific thresholds from the validation set.

    \item \textbf{The choice of layer matters as much as the choice of method.}
    At \texttt{layer4} GradCAM and LayerCAM are almost indistinguishable for most classes.
    The real difference emerges at \texttt{layer3}: LayerCAM's element-wise weighting preserves fine spatial detail that GradCAM's global pooling discards.
    Comparing CAMs at multiple layers is cheap and highly informative.

    \item \textbf{FP and TN cases are the most informative for explanation analysis.}
    True positives often look plausible almost by definition, but false positives and true negatives reveal much more about how the model actually behaves. In low-confidence and TN cases, GradCAM often shifted toward image corners and borders, suggesting possible reliance on artifacts or framing effects, while LayerCAM more often stayed within the chest area. Even when LayerCAM was not clearly correct, its maps were usually more anatomically plausible, which made such cases especially useful for comparing the practical behaviour of the two methods.

    \item \textbf{Domain knowledge cannot be replaced by AUC alone.}
    Evaluating whether a highlighted region is anatomically meaningful requires expertise we do not have.
    In the absence of radiologist annotations, LLM-assisted validation (with appropriate scepticism) was the most practical option available to us.
\end{itemize}

% ─────────────────────────────────────────────────────────────────────
\section{Further Goals}

\begin{itemize}
    \item \textbf{Train on more data.}
    The biggest limitation we have right now is overfitting --- our training AUC reached 0.9999 while validation AUC stayed at 0.83. The main reason is that we only used 10k images out of 224k available, because training already took over an hour on our hardware. In the future, we would like to train on the full dataset and also try data augmentation and stronger regularisation to help the model generalise better.

    \item \textbf{Validate explanations with a real expert.}
    Right now we used LLM assistance to judge whether the highlighted regions make anatomical sense, which is a reasonable workaround but not a proper evaluation. Ideally, we would want a radiologist to look at the heatmaps and tell us whether GradCAM and LayerCAM are actually focusing on the right areas. Without that, we cannot really say how clinically useful the explanations are.

    \item \textbf{Compare with other XAI methods.}
    We only implemented GradCAM and LayerCAM, but there are other explanation methods worth trying --- for example ScoreCAM, LIME, or SHAP. It would be interesting to see how they behave on the same images and whether they agree or disagree with our heatmaps, especially in the difficult FP and TN cases.
\end{itemize}

% ─────────────────────────────────────────────────────────────────────
\section*{References}

\begin{itemize}[leftmargin=*]
    \item Jiang et al., \textit{LayerCAM: Exploring Hierarchical Class Activation Maps for Localization}, 2021.
    \item Selvaraju et al., \textit{Grad-CAM: Visual Explanations from Deep Networks via Gradient-based Localization}, ICCV 2017.
    \item Irvin et al., \textit{CheXpert: A Large Chest Radiograph Dataset with Uncertainty Labels and Expert Comparison}, AAAI 2019.
    \item He et al., \textit{Deep Residual Learning for Image Recognition}, CVPR 2016.
\end{itemize}

\end{document}
